\chapter{ディレクトリとパス}

\section*{ファイル・システムの全体像}

 UNIX システムでは、ファイル・システムは、ツリー状の階層化ディレクトリ構造を
している。
これは簡単に言うと次のようになる。
\begin{enumerate}
\item たくさんあるファイルを、種類ごとに\textbf{「ディレクトリ」}と
呼ばれる箱に分けてしまうことで整理する。
\item \textbf{ディレクトリ}の中にも、また\textbf{ディレクトリ}を作れ、
更に分類して整理できる。
\item ディレクトリの構造図を書くと、木が枝をひろげたよう(ツリー状)に見える。
\end{enumerate}

このようなディレクトリ構造は、後に UNIX システムを真似した MS-DOS や
 Windows などにも採り入れられた。
 Windows では、ディレクトリのことをフォルダと呼んでいる。

\bigskip

 UNIX システムには、いろいろなファイルやコマンドがある。
ディレクトリの中を自由に移動できるようになれば、
これらのことを詳しく知ることができるようになるだろう。

また、自分のファイルをディレクトリに分けてしまっておくと、
きれいに整理されて、とても便利に使える。

この章では、ディレクトリの移動、作成、消去などについて紹介しよう。

\section{UNIX システムのディレクトリ構造}

\subsection{ディレクトリの全体像}

最初に UNIX システムの全体像を紹介してしまう。

次のページの図 \ref{dirstru} は、 UNIX のディレクトリ構造の主な部分を書いた
ものである。
これには、次のようなディレクトリがある。

\begin{itemize}
\setlength{\parindent}{1zw}
\item \textbf{\large ルート・ディレクトリ /}

一番上のディレクトリを「\textbf{ルート・ディレクトリ}」といい、
「\texttt{/}」で
表す。
なお、「ルート -- root 」というのは「(木の)根」のことだ。

ルートの下には、「\texttt{home}」、「\texttt{bin}」、「\texttt{etc}」など、
いろいろなディレクトリが、枝をひろげている。
そして、それらのディレクトリの先にも、またディレクトリがあり、
さらに枝がひろがっている。

これらがツリー状に見える。
絵としては、木が上下さかさまに描かれていると思ってくれればいい。

\begin{figure}[H]
\centering{\includegraphics[scale=0.8]{dirstru.eps}}
\caption{UNIX システムのディレクトリ構造}\label{dirstru}
\end{figure}


では、ルート・ディレクトリの下を一つずつ説明していこう。

\item \textbf{\large ホーム・ディレクトリ home}

まず、「\texttt{/}」の下には、「\texttt{home}」というディレクトリがあり、
この中に人の名前のついたディレクトリがある。
これらは、 UNIX システムで \textbf{「ホーム・ディレクトリ」}と呼ばれている。

たとえば上の図では、 \texttt{kokubo} さんのホーム・ディレクトリは、
「\texttt{/}」の下の「\texttt{home}」の下の「\texttt{kokubo}」だ。
そしてこれを、「\texttt{/home/kokubo}」と表す。
ディレクトリの区切りを表すマークは「\texttt{/}」である。

また同様に、 \texttt{tomoda} さんのホーム・ディレクトリは、「\texttt{/}」
の下の「\texttt{home}」の下の「\texttt{tomoda}」で、
これを「\texttt{/home/tomoda}」と表す。

 login した直後、それぞれのユーザは、自分のホーム・ディレクトリにまず入る。
ここまでの演習で作ったファイルも、みんなホーム・ディレクトリにあるはずだ。

上の例では、 \texttt{kokubo} さんは login すると、
「\texttt{/home/kokubo}」に入る。
そして、ここまでの演習で \texttt{kokubo} さんの作ったファイルは、
ここに置いてあるということだ。

\item \textbf{\large bin}

「\texttt{bin}」というのは「バイナリ -- binary」の略で、本来は、
文字ではないデータが入っているファイルのことだ。
ちなみに、マシン語で書かれたコマンドや、画像、音声データなどが、
本来の意味ではバイナリである。
なお、 UNIX システムでは、コマンドを入れるためのディレクトリに
「\texttt{bin}」という名前を付けている。

「\texttt{/}」の下の「\texttt{bin}」、つまり「\texttt{/bin}」には、
 UNIX システムの一番基本的なコマンドが置いてある。

\item \textbf{\large etc}

「\texttt{etc}」は「エトセトラ -- etc.」である。
 UNIX システムでは、ここに各種システム設定用のファイルが置かれている。

\item \textbf{\large tmp}

「\texttt{tmp}」は「テンポラリ -- temporary」である。
主に、システムが一時的にファイルを作るときに、ここを使う。

\item \textbf{\large usr}

「\texttt{usr}」は「ユーザ -- user」という意味だ。

この中は更に分かれていて、「\texttt{/usr/bin}」にはユーザが使う各種コマンド、
「\texttt{/usr/include}」には C 言語のヘッダー・ファイルというものが、
「\texttt{/usr/lib}」にはユーザ・コマンドが使うライブラリというデータが
置かれている。

また、「\texttt{/usr/local}」以下には、オプションのコマンドなどが
置かれている。
オプションのコマンドは、 UNIX の基本システムには付いて来ないものだが、
便利なものが多い。
例えば、 Mule などもオプションのコマンドである。

なお、この図には書いていないが、「\texttt{/usr/local}」の下にも、
更に「\texttt{bin}」や「\texttt{lib}」などがある。

\item \textbf{\large var}

「\texttt{var}」は「変化する -- variable」という意味。
 login の記録や、メールや、システムの記録などの、
日々変化するものが置かれている。
\end{itemize}

\newpage

\subsection{B 演習室の実際}

ここまでで紹介してきたのは、あくまでも標準のスタイルだ。
ディレクトリの構造は、システムの管理者の好みに応じて
変えることができる。

B 演習室では、ホーム・ディレクトリの様子は、図 \ref{Benshu} のように
変更されている。
なお、ホーム・ディレクトリ以外は標準の構成のままである。

図を見てもらうとわかるが、「\texttt{/}」の下に「\texttt{home}」があるところ
までは標準と一緒だ。
違うのは、この中が更に学年によって「\texttt{inf8}」、
「\texttt{ei97}」、「\texttt{ei98}」、「\texttt{ei99}」と切り分けてあるところだ。
そして、その中にそれぞれのユーザのホーム・ディレクトリがある。

たとえば、 \texttt{ei99001} さんのホーム・ディレクトリは
「\texttt{/home/ei99/ei99001}」だ。
同様に \texttt{ei99002} さんは「\texttt{/home/ei99/ei99002}」となる。

そして、この章の後の方で説明するが、ユーザは
自分でホーム・ディレクトリの中に「\texttt{Report}」だとか、
 \texttt{Sono1} とかのディレクトリを作って、ファイルを整理することが
できる。

\begin{figure}[H]
\centering{\includegraphics[scale=0.8]{Benshu.eps}}
\caption{B 演習室の実際のホーム・ディレクトリの構造}\label{Benshu}
\end{figure}

\section{パス}

 UNIX システムでは、
ファイルやディレクトリの場所を指定するのに、
「\textbf{パス}」を使う。

「パス -- path」というのは、通り道や、経路という意味だ。
パスには、「\textbf{絶対パス}」と「\textbf{相対パス}」がある。

このパスをマスターすることで、好きな場所に移動して UNIX の中を探検したり、
ファイルを好きな場所にコピーしたりすることができる\footnote{セキュリティ上、
いじれないようになっている場所もある。}ようになる。

パスは UNIX システムをマスターする上で、最重要ポイントの一つだ。

\subsection{根っこからの通り道: 絶対パス}

既に、ここまでのところで紹介してしまったが、
「\texttt{/home/kokubo}」のように、 \underline{一番上のルート・}\\
\underline{ディレクトリから、順番にディレクトリを書いて表す}
方法を、\textbf{「絶対パス」}と呼ぶ。

\bigskip

図 \ref{absol} を使って説明しよう。
まず、一番上にルート・ディレクトリがあり、
これは絶対パスで表すと「\texttt{/}」になる。

「\texttt{/}」の下には「\texttt{home}」があって、
これは絶対パスでは「\texttt{/home}」となる。


その下は学年ごとにディレクトリが用意されていて、「\texttt{ei99}」は
絶対パスでは「\texttt{/home/ei99}」となる。

たとえばその中に、 \texttt{ei99000} さんのホーム・ディレクトリがあって、
これは「\texttt{/home/ei99/ei99000}」だ。

そして、 \texttt{ei99000} さんは、自分のホーム・ディレクトリの中に
「\texttt{Report}」というディレクトリを作ることができる。
これは「\texttt{/home/ei99/ei99000/Report}」と表される。

またその中にディレクトリやファイルを作ることもできる。
たとえば「\texttt{11gatsu}」というディレクトリを作ると、
これは絶対パスでは「\texttt{/home/ei99/ei99000/Report/11gatsu}」になる。

このように根っこ(ルート・ディレクトリ)から順番に、書いていくのが絶対パスである。

\bigskip

なお、ここまでディレクトリだけしか説明しなかったが、ファイルの場合も
全く同様である。
 \texttt{ei99000} さんのホーム・ディレクトリ「\texttt{/home/ei99/ei99000}」に
「\texttt{ThisMonth}」というファイルがある場合、
このファイルは絶対パスでは「\texttt{/home/ei99/ei99000/ThisMonth}」と書く。

\begin{figure}[H]
\centering{\includegraphics[scale=0.8]{abosule.eps}}
\caption{絶対パスでの表し方}\label{absol}
\end{figure}

\subsection{自分がいるところが中心: 相対パス}

絶対パスを使って、ルート・ディレクトリから、順番書いていけば、
どんなディレクトリでも表せる。
しかし、絶対パスだと、だんだんディレクトリが下の方になってくると、
どんどん長くなって面倒だ。

\bigskip

そこで、 UNIX システムでは、\textbf{「相対パス」}という、
もう一つの方法が用意されている。
相対パスでは、俺が世界の中心という感じで、自分が今いるところを基準にして
指定する。

\bigskip

では、図 \ref{relative} を見ながら説明しよう。
これは、さっき絶対パスの説明に使ったものと全く同じだ。

まず、 \texttt{ei99000} さんは、今自分のホーム・ディレクトリにいるとしよう。
この下に、「\texttt{Report}」というディレクトリがあるとする。

するとこれは、相対パスでは、単に「\texttt{Report}」と書く。

絶対パスだと、「\texttt{/home/ei99/ei99000/Report}」と長くて面倒だが、
相対パスではこんなに短くなる。

それから、「\texttt{Report}」の下の「\texttt{11gatsu}」も、
相対パスでは「\texttt{Report/11gatsu}」と、とても簡単に書ける。

このように相対パスとは、
\underline{自分が今いるところから、順番にパスを書いていく}方法である。

\begin{figure}[H]
\centering{\includegraphics[scale=0.8]{relative1.eps}}
\caption{相対パスので表し方}\label{relative}
\end{figure}

\subsubsection{1つ上のディレクトリ: 「\texttt{..}」}

ここまでは自分よりも下のディレクトリの表し方を紹介してきた。
では、逆に上の方はどうなるのかを紹介しよう。

 UNIX システムでは、 1 つ上のディレクトリを「\texttt{..}」と表す。

つまり、自分のホームの一つ上の「\texttt{/home/ei99}」は、
「\texttt{..}」と表される。

この\textbf{「ピリオド 2 つで、 1 つ上」}は、とてもよく使うので覚えておこう。

\bigskip

では、 2 つ上はどうか？

1 つ上が「\texttt{..}」で、ディレクトリの区切りが「\texttt{/}」なので、
2 つ上は「\texttt{../..}」となる。

つまり、「\texttt{/home}」は、相対パスでは「\texttt{../..}」となる。

同様にルート・ディレクトリ「\texttt{/}」は、相対パスでは
「\texttt{../../..}」である。

\subsubsection{自分が今いるディレクトリ: 「\texttt{.}」}

最後にもう一つだけ紹介しておく。
\textbf{自分が今いる場所は、相対パスでは「\texttt{.}」}と、
ピリオド 1 つで表わす。

これもよく使うので覚えておこう。

たとえば、今「\texttt{/home/ei99/ei99000}」にいるなら、ここは相対パスで
表すと「\texttt{.}」となる。

\subsubsection{自分のいる場所によって相対パスは変わる}

なお、以上の例は、あくまで「\texttt{/home/ei99/ei99000}」にいる場合の話だ。

相対パスは自分が今いるところを中心に指定するので、
別なディレクトリにいる場合、話が全く変わってくる。

たとえば、「\texttt{/home/ei99/ei99000/Report}」にいるとしよう。

今度は、自分のホーム・ディレクトリは一つ上にあるので、
相対パスでは「\texttt{..}」となる。

また、「\texttt{11gatsu}」は、今いるところのすぐ下にあるので、相対パスでは
「\texttt{11gatsu}」になる。

このように相対パスは、現在どこのディレクトリに自分がいるかによって
変わってしまう。

\begin{figure}[H]
\centering{\includegraphics[scale=0.8]{relative2.eps}}
\caption{\texttt{Report} の中にいるときの相対パス}\label{relative2}
\end{figure}

\subsection{絶対パスと相対パス、どっちがお特?}

 UNIX システムでは、絶対パスで指定しても、相対パスで指定しても、
効果は全く同じである。

だから、自分で好きな方を使っていい。

ただ、自分が今いるところの近くを指定するには「相対パス」を、
ルート・ディレクトリの近くを指定するには「絶対パス」を使った方が簡単だ。

このように使い分けるので、両方を知っておいた方が便利だ。

\bigskip
\bigskip

\begin{BOXNOTE}
\textbf{\large ☆練習☆}
\begin{enumerate}
\item 下の図の \fbox{1} の絶対パスは？
\item では、今、あなたは \fbox{a} にいる。すると \fbox{1} の相対パスは？
\item では、あなたが \fbox{b} にいるときの、 \fbox{1} の相対パスは？
\item \fbox{2} の絶対パスは？
\item では、今あなたが \fbox{a} にいるとき、 \fbox{2} の相対パスは？
\item では、今あなたが \fbox{b} にいるときはの \fbox{2} の相対パスは？
\end{enumerate}
\begin{figure}[H]
\centering{\includegraphics[scale=0.8]{quiz.eps}}
\end{figure}

\bigskip
\textbf{\large 答え}

\begin{enumerate}
\item ~~~\texttt{/home/ei99}
\item ~~~\texttt{..}
\item ~~~\texttt{ei99}
\item ~~~\texttt{/home/ei99/ei99069}
\item ~~~\texttt{../ei99069}
\item ~~~\texttt{ei99/ei99069}
\end{enumerate}
\end{BOXNOTE}
\newpage

\section{ディレクトリ操作コマンド}

では、具体的に、コマンドを使ってディレクトリを移動したり、
ディレクトリを作ったりしてみよう。

\subsection{現在いるディレクトリは?: \texttt{pwd}}

login した直後にいる場所が、きみのホーム・ディレクトリだと、前の方で説明した。
では、それがどこか見てみよう。

今いるディレクトリを表示するには、 \texttt{pwd} コマンドを使う。
 \texttt{pwd} は「\textbf{p}rint \textbf{w}orking \textbf{d}irectory」の略で、
「今お仕事をしているディレクトリを表示してね」という意味だ。

使い方は、単に次のように打てばいい。

\begin{center}
\begin{minipage}{1.8cm}
\begin{shadebox}
\smallskip
{\Large \texttt{pwd}}
\smallskip
\end{shadebox}
\end{minipage}
\end{center}

さて今、図 \ref{now} のようなディレクトリ構造で、
\texttt{ei99000} さんは、「\texttt{/home/ei99/ei99000}」にいると
しよう。
このとき、 \texttt{pwd} を実行すると、

\begin{progls}
% \slbfl{pwd}
/home/ei99/ei99000
%
\end{progls}

となり、 今いるのが「\texttt{/home/ei99/ei99000}」だと表示される。

\bigskip

\begin{figure}[H]
\centering{\includegraphics[scale=0.8]{abrel.eps}}
\caption{\texttt{ei99000} さんが「\texttt{/home/ei99/ei99000}」にいるときの
絶対パスと相対パス}\label{now}
\end{figure}

\newpage

\subsection{ディレクトリを移動: \texttt{cd}}

では、ディレクトリの冒険の旅に出かけよう。
ディレクトリを移動するには、 \texttt{cd} を使う。
これは「\textbf{c}hange \textbf{d}irectory -- チェンジ・ディレクトリ」の略だ。
使い方は、
\begin{center}
\begin{minipage}{4cm}
\begin{shadebox}
\smallskip
{\Large \texttt{cd} 「パス」}
\smallskip
\end{shadebox}
\end{minipage}
\end{center}

\noindent
である。
この「パス」というのは、絶対パス、相対パスのどちらでもいい。

\subsubsection{ 1 つ上へ移動}

では、 1 つ上のディレクトリに移動してみよう。
現在、図 \ref{now} のように、「\texttt{/home/ei99/ei99000}」としよう。

 1 つ上のディレクトリは、絶対パスでは「\texttt{/home/ei99}」で、
相対パスでは「\texttt{..}」だ。

よって、 1 つ上に移動するには、「\texttt{cd ..}」と、
「\texttt{cd /home/ei99}」のどちらでもいい。

とは言え、明らかに「\texttt{cd ..}」の方が楽なので、こちらを実行しよう。

\begin{progls}
% \slbfl{cd ..}
%
\end{progls}

これで、ディレクトリを一つ登れた。
 \texttt{pwd} で確認してみよう。

\begin{progls}
% \slbfl{pwd}
/home/ei99
%
\end{progls}

確かに、「\texttt{/home/ei99}」にいることがわかる。

\bigskip

では、ここにどんなファイルがあるか \texttt{ls} で見てみよう。
すると、情報システム工学科 1 年みんなのディレクトリが見える。

\begin{progls}
% \slbfl{ls}
ei99001 ei99006 ei99011 ei99016 ei99021 ei99026 ei99031 ei99036 ei99041
ei99002 ei99007 ei99012 ei99017 ei99022 ei99027 ei99032 ei99037 ei99042
ei99003 ei99008 ei99013 ei99018 ei99023 ei99028 ei99033 ei99038 ei99043
ei99004 ei99009 ei99014 ei99019 ei99024 ei99029 ei99034 ei99039 ei99044
ei99005 ei99010 ei99015 ei99020 ei99025 ei99030 ei99035 ei99040 ei99100
%
\end{progls}

ここで、 \texttt{ls -F} と打ってみよう。

\texttt{-F} は「\textbf{f}ile -- ファイル(の種類)」の意味で、
次のようにディレクトリの後には \texttt{/} を付けて表示してくれるので、
ずいぶん見やすくなる。

\begin{progls}
% \slbfl{ls -F}
ei99001/        ei99010/        ei99019/        ei99028/        ei99037/
ei99002/        ei99011/        ei99020/        ei99029/        ei99038/
ei99003/        ei99012/        ei99021/        ei99030/        ei99039/
ei99004/        ei99013/        ei99022/        ei99031/        ei99040/
ei99005/        ei99014/        ei99023/        ei99032/        ei99041/
ei99006/        ei99015/        ei99024/        ei99033/        ei99042/
ei99007/        ei99016/        ei99025/        ei99034/        ei99043/
ei99008/        ei99017/        ei99026/        ei99035/        ei99044/
ei99009/        ei99018/        ei99027/        ei99036/        ei99100/
%
\end{progls}

 \texttt{ls -F} の \texttt{-F} のように、「\texttt{-なんとか}」のような
形の引数を\textbf{「コマンドのオプション」}という。

コマンドには、それぞれのコマンドによって、いろいろなオプションがある。
そして、オプションを付けると、コマンドの働きを変えることができるようになっている。
どんなオプションがあるかはオンライン・マニュアルに載っている。

\subsubsection{一気に絶対パスで移動}

次は、絶対パスを使って、一気に「\texttt{/bin}」に移動してみよう。

「\texttt{/bin}」は、 UNIX の最も基本的なコマンドの置いてあるディレクトリである。

では「\texttt{cd /bin}」と打とう。

\begin{progls}
% \slbfl{cd /bin}
%
\end{progls}

これで、一気に「\texttt{/bin}」まで移動できたはずだ。
 \texttt{pwd} で確認してみよう。

\begin{progls}
% \slbfl{pwd}
/bin
%
\end{progls}

確かに「\texttt{/bin}」にいることがわかる。

では、どんなファイルがあるか \texttt{ls} で見てみよう。

\begin{progls}
% \slbfl{ls}
[               dd              kill            pwd             sleep
cat             df              ln              rcp             stty
chio            domainname      ls              red             sync
chmod           echo            mkdir           rm              test
cp              ed              mv              rmail
csh             expr            pax             rmdir
date            hostname        ps              sh
%
\end{progls}

なお、ここで \texttt{ls -F} すると、次のようになる。
ここで、 \texttt{*} マークは実行できるコマンドの目印である。

\begin{progls}
% \slbfl{ls -F}
[*              dd*             kill*           pwd*            sleep*
cat*            df*             ln*             rcp*            stty*
chio*           domainname*     ls*             red*            sync*
chmod*          echo*           mkdir*          rm*             test*
cp*             ed*             mv*             rmail*
csh*            expr*           pax*            rmdir*
date*           hostname*       ps*             sh*
%
\end{progls}

「\texttt{cd}」は、ここまで見てきたように、「相対パス」と「絶対パス」の両方が
使える。

じゃ、どっち使えばいいのかというと、
「自分のいるとこの近くなら、相対パス」、「遠くなら絶対パス」が
便利なことが多い。

\subsubsection{ホームに戻る}

では、ホーム・ディレクトリに戻ることにしよう。

これはとても簡単で、単に「\texttt{cd}」と打てばいい。
いつ、どこにいても、自分のホーム・ディレクトリには、「\texttt{cd}」と
打つだけで戻れる。

では、やってみよう。

\begin{progls}
% \slbfl{cd}
%
\end{progls}

\texttt{pwd} で確認してみよう。

\begin{progls}
% \slbfl{pwd}
/home/ei99/ei99000
%
\end{progls}

確かに、ホーム・ディレクトリに戻れたことがわかる。

\newpage

\subsection{ディレクトリを作る: \texttt{mkdir}}

ディレクトリを作るには \texttt{mkdir} だ。
これは「\textbf{m}ake \textbf{dir}ectory -- メイク・ディレクトリ」で、
そのままの意味だ。

使い方は次の通り。

\begin{center}
\begin{minipage}{10cm}
\begin{shadebox}
\smallskip
{\Large\bf \texttt{mkdir} 「作りたいディレクトリの名前」}
\smallskip
\end{shadebox}
\end{minipage}
\end{center}


では、自分のホームの下に「\texttt{Enshu}」とかいうディレクトリを作ってみよう。

\begin{progls}
% \slbfl{mkdir Enshu}
%
\end{progls}

では、 \texttt{ls} で確認してみよう。

\begin{progls}
% \slbfl{ls}
1998nen         2000nen         RMAIL           ThisYear       exam1
1999nen         Enshu           ThisMonth       Years
%
\end{progls}

これでは、「\texttt{Enshu}」があるのはわかるが、ファイルかディレクトリか
区別がつかない。

区別をはっきりつけるには \texttt{ls -F} と打つ。

\begin{progls}
% \slbfl{ls}
1998nen         2000nen         RMAIL           ThisYear       exam1
1999nen         Enshu/          ThisMonth       Years
%
\end{progls}

たしかに、「\texttt{Enshu}」というディレクトリがあるのがわかる。

では、「\texttt{Enshu}」の中に入ってみよう。
これは、絶対パスなら「\texttt{/home/ei99/ei99000/Enshu}」で、
相対パスでは「\texttt{Enshu}」なので、相対パスの方が簡単だ。
よって、「\texttt{cd Enshu}」と打とう。

\begin{progls}
% \slbfl{cd Enshu}
%
\end{progls}

では、 \texttt{pwd} で確認してみよう。

\begin{progls}
% \slbfl{pwd}
/home/ei99/ei99000/Enshu
%
\end{progls}

\subsection{ディレクトリを消す: \texttt{rmdir}}

ディレクトリを消すには、 \texttt{rmdir} だ。
これは「\textbf{r}e\textbf{m}ove \textbf{dir}ectory -- ディレクトリ削除」の
略だ。
使い方は次の通りだ。

\begin{center}
\begin{minipage}{10cm}
\begin{shadebox}
\smallskip
{\Large\bf \texttt{rmdir} 「消したいディレクトリの名前」}
\smallskip
\end{shadebox}
\end{minipage}
\end{center}

では、実際にやってみよう。
まず、自分のホーム・ディレクトリにもどろう。

ホームに戻るには単に \texttt{cd} でいい。

\begin{progls}
% \slbfl{cd}
%
\end{progls}

次のように \texttt{pwd} で確認すると、確かにホームにいることがわかる。

\begin{progls}
% \slbfl{pwd}
/home/ei99/ei99000
%
\end{progls}

では、次のようにして、さっき作った「\texttt{Enshu}」を消してみる。

\begin{progls}
% \slbfl{rmdir Enshu}
%
\end{progls}

これで、消えたはずだ。
\texttt{ls -F} で確認してみよう。

\begin{progls}
% \slbfl{ls -F}
1998nen         2000nen         ThisMonth       Years
1999nen         RMAIL           ThisYear        exam1
%
\end{progls}

確かに消えている。

\bigskip

注・ディレクトリの中に何かファイルが残っていると、 \texttt{rmdir} で
ディレクトリを消すことはできない。

\subsection{ファイルの移動、コピー}

\texttt{mv} や \texttt{cp} で別のディレクトリにファイルを移動したり、
コピーする方法を紹介しよう。

たとえば、「\texttt{CAL}」というディレクトリを新しく作って、
「\texttt{ThisMonth}」というファイルをその中に移動するには次のようにする。

まず、「\texttt{CAL}」というディレクトリがあるかどうか確認する。

\begin{progls}
% \slbfl{ls -F}
1998nen         2000nen         ThisMonth       Years
1999nen         RMAIL           ThisYear        exam1
%
\end{progls}

すると、上のようにないことがわかる。
そこで「\texttt{CAL}」を \texttt{mkdir} で作る。

\begin{progls}
% \slbfl{mkdir CAL}
%
\end{progls}

ファイルをディレクトリの中に移動やコピーする場合でも、
 \texttt{mv} や \texttt{cp} の使い方は、普通に移動したりコピーするときと
全く同じだ。
ただ、ファイル名として、パスを書いてやればいい。

つまり、 \texttt{ThisMonth} を \texttt{CAL} に移動するには、
「\texttt{mv ThisMonth CAL/ThisMonth}」となる。

ちなみに、これは「\texttt{mv ThisMonth CAL/.}」とも書ける。

「\texttt{.}」をファイル名の代わりに使うと、「同じ名前で」という
意味になるのだ。

このように「\texttt{.}」は、ディレクトリのときには「現在の場所」、
ファイル名のときには「同じ名前で」という意味になるので、混同しないように
注意しよう。

\bigskip

では、実際にやってみる。

\begin{progls}
% \slbfl{mv ThisMonth CAL/.}
%
\end{progls}

\texttt{ls -F} で確かめてみると、たしかに「\texttt{ThisMonth}」がないことがわかる。
\begin{progls}
% \slbfl{ls -F}
1998nen         2000nen         RMAIL           Years
1999nen         CAL/            ThisYear        exam1
%
\end{progls}

「\texttt{CAL}」に入って中を見て見ると、「\texttt{ThisMonth}」があることがわかる。
\begin{progls}
% \slbfl{cd CAL}
% \slbfl{ls -F}
ThisMonth
%
\end{progls}

なお、この説明では \texttt{mv} を使ってファイルを移動したが、代わりに \texttt{cp} を使えば、
コピーになる。

\section*{この章で紹介したコマンド}

\begin{center}
\begin{minipage}{13cm}
\begin{itembox}[l]{ディレクトリ操作コマンド}
\begin{description}
\item[\texttt{pwd} :] 自分が現在いる場所を表示

使い方: \texttt{pwd}
\item[\texttt{mkdir} :] ディレクトリの作成

使い方: \texttt{mkdir} 「ディレクトリ名」
\item[\texttt{rkdir} :] ディレクトリの削除

使い方: \texttt{rmdir} 「ディレクトリ名」
\end{description}
\end{itembox}
\end{minipage}
\end{center}
